\documentclass[12pt]{article}

\usepackage{url,verbatim, amsmath, amssymb, amsfonts, dsfont, color}
\textwidth 15true cm
\textheight 9true in
\oddsidemargin 0.30true in
\evensidemargin 0.30true in
\topmargin -0.3true in
\headsep 0.4true in

\def\sgn{{\rm sign}}
\def\drightarrow{{\stackrel{d}{\rightarrow}}}
\def\dotsim{\stackrel{.}{\sim}}
\def\goesd{\stackrel{d}{\rightarrow}}
\def\goesp{\stackrel{p}{\rightarrow}}
\def\goeslaw{\stackrel{{\cal L}}{\longrightarrow}}
%
\def\vp{{\varphi}}
\def\E{{\rm E}}
\def\Pvalue{{$p$-value}}
\def\Pvalues{{$p$-values}}
\def\pr{{\rm pr}}
\def\Pr{{\rm pr}}
\def\logit{{\rm logit}}
\def\T{{ \mathrm{\scriptscriptstyle T} }}
\def\diag{{\rm diag}}  
\def\half{{\textstyle{1\over 2}}}
\def\ith{{$i{\rm th}$ }}
\def\var{{\rm var}}
\def\cl{{{\rm c}\ell}}
\def\checkx{{\checked\hspace{-6pt}\setminus}}
\def\grad{{\nabla}}
\def\bs#1{{\boldsymbol{#1}}}

\newcommand{\R}{\mathbb{R}}
\newcommand{\N}{\mathbb N}
\newcommand{\Z}{\mathbb Z}
\newcommand{\Q}{\mathbb Q}
\newcommand{\1}{\mathds 1}
\newcommand{\Var}{\text{Var}}
\newcommand{\Cov}{\text{Cov}}
\newcommand{\sd}{\text{sd}}
\newcommand{\se}{\text{se}}
\renewcommand{\L}{\mathcal {L}}
\renewcommand{\lambda}{\uplambda}
\newcommand{\iidsim}{\stackrel{iid}{\sim}}
\newcommand{\otherwise}{\text{otherwise}}
\newcommand{\indep}{\perp \!\!\! \perp}
\renewcommand{\epsilon}{\varepsilon}
\newcommand{\st}{\text{ s.t. }}
\newcommand{\ds}{\displaystyle}
\newcommand{\ind}{\stackrel{d}{\to}}
\newcommand{\inp}{\stackrel{p}{\to}}
\newcommand{\inas}{\stackrel{a.s.}{\to}}
\DeclareMathOperator*{\argmax}{arg\,max}
\DeclareMathOperator*{\argmin}{arg\,min}

\sloppy %permits line-breaking of urls

\def\blue{\color{blue}}
\def\red{\color{red}}
\def\sgn{{\rm sign}}
\def\drightarrow{{\stackrel{d}{\rightarrow}}}
\begin{document}

\noindent{\bf Questions Week 6} {\hfill STA 2212S 2026}\\

%{\bf Due January 14} \hfill{\it MS, Exercise 5.2}

\bigskip

\begin{enumerate}

\item {\it SM Exercise 7.3.8} If $I$ is Bernoulli with $\Pr(I=1) = (\alpha_2-\alpha)/(\alpha_2-\alpha_1)$, and ${\cal R}_{\alpha_1}, {\cal R}_{\alpha_2}$ are rejection regions of size $\alpha_1, \alpha_2$, show that the rejection region $${\cal R} = I{\cal R}_{\alpha_1} + (1-I){\cal R}_{\alpha_2}$$ has size $\alpha$. {\it Note: this is the basis for randomized tests for discrete distributions.}

\textcolor{red}{
We want to show that $P(X \in \mathcal R) = \alpha$. This rejection rule implies a test function $\phi(X, I)$ defined as 
\begin{equation*}
    \phi(X, I) = I\1_{X \in \mathcal R_{\alpha_1}} + (1 - I) \1_{X \in \mathcal R_{\alpha_2}}
\end{equation*}
Then 
\begin{equation*}
    P_{H_0} (X \in \mathcal R) = P_{H_0} (\phi(X, I) = 1) 
\end{equation*}
Since $\phi(X, I)$ is binary, then 
\begin{align*}
    P_{H_0}(\phi(X, I) = 1) &= E_{H_0} (\phi(X, I))\\
    &= P(I = 1)P(X \in \mathcal R_{\alpha_1}) + P(I = 0) P(X \in \mathcal R_{\alpha_2})\\
    &= \frac{\alpha_2 - \alpha}{\alpha_2 - \alpha_1} \alpha_1 + \frac{\alpha - \alpha_1}{\alpha_2 - \alpha_1} \alpha_2\\
    &= \frac{\alpha_1\alpha_2 - \alpha_1\alpha + \alpha_2\alpha -\alpha_1\alpha_2}{\alpha_2 - \alpha_1}\\
    &= \alpha
\end{align*}
hence $\mathcal R$ has size $\alpha$. 
}

\item {\em  A Bayesian approach to  multiple testing} (adapted from Efron and Hastie \S 15.3)\\
Suppose we have computed $m$ test statistics $t_1, \dots, t_m$, and for each $j$ we will reject $H_{0j}$ if $t_j > t_0$, where $t_0$ is some cut-off to be determined. A statistical model for this could start with prior probabilities $\pi_0$ and $\pi_1 = 1-\pi_0$ that any $H_{0j}$ is true or not; and densities for $t_j$ either $f_0(t)$ or $f_1(t)$,  when $H_{0j}$ is true, or false, respectively. Since we don't know which $H_{0j}$ are true, the marginal density of $t_j$ is
\begin{equation}
f(t) = \pi_0 f_0(t) + \pi_1f_1(t).
\end{equation}
\begin{description}
\item{(i)} The {\em Bayes false-discovery rate} for cutoff $t_0$ is defined as
$$
Fdr(t_0) = \pr\{ H_{0j} \text{ is true } \mid t_j \ge t_0\}.
$$
Show that this is equal to $$\pi_0\{1-F_0(t_0)\}/\{1-F(t_0)\},$$ 
where $F(t)$ is the cumulative distribution function for the density (1). 
\item{(ii)} We don't know $F(t_0)$, but a reasonable estimate is the empirical c.d.f   $\widehat F(t_0) = m^{-1}\sum_{j=1}^m 1\{t_j \le t_0\}$ and hence
$$
\widehat{Fdr}(t_0)=\pi_0\{1-F_0(t_0)\}/\{1-\widehat F(t_0)\},
$$  Show that for $0 < q < 1$, the condition $$p_{(i)}\le (i/m)q$$ is equivalent to $$\widehat{Fdr}(t_{(i)}) \le \pi_0q.$$ 
\end{description}

\textcolor{red}{
(i) By Bayes' rule, 
\begin{align*}
    P(H_{0j} \text{ is true} \mid t_j \geq t_0) &= \frac{P(H_{0j} \text{ is true})}{P(t_j \geq t_0)} P(t_j \geq t_0 \mid H_{0j} \text{ is true})\\
    &= \frac{\pi_0}{1 - F(t_0)} (1 - F_0 (t_0))
\end{align*}
since the density of $t_j$ given $H_{0j}$ is true is $f_0$ and the marginal density of $t_j$ is $f$. 
}

\textcolor{red}{
(ii) Consider the $p$-values $p_j = 1 - F_0(t_j)$, so the ordering $p_{(1)} \leq p_{(2)} \leq \cdots \leq p_{(m)}$ of the $p$-values corresponding to the ordering $t_{(1)} \geq t_{(2)} \geq \cdots \geq t_{(m)}$ of the test statistics. Then 
\begin{align*}
    &p_{(i)} \leq \frac im q \\
    \iff & 1 - F_0(t_{(i)}) \leq \frac im q\\
    \iff & \pi_0 \frac{1 - F_0(t_{(i)})}{1 - \hat F(t_{(i)})} \leq \frac im q \frac{\pi_0}{1 - \hat F(t_{(i)})}
\end{align*}
By definition, 
\begin{equation*}
    \hat F(t_{(i)}) = \frac 1m \sum_{j=1}^m \1_{t_j \leq t_{(i)}} = \frac{m-i}{m}
\end{equation*}
since $t_{(i)}$ is the $i$-th smallest test statistic. 
Then this implies $1 - \hat F(t_{(i)}) = i/m$, hence 
\begin{equation*}
    p_{(i)} \leq \frac im q \iff \widehat{Fdr}(t_{(i)}) \leq \pi_0 q 
\end{equation*}
}

\item State and prove the Neyman-Pearson Lemma. 

\textcolor{red}{
\textbf{Lemma:} If we consider parameter space $\Theta = \{\theta_0, \theta_1\}$ and a test $H_0: \theta = \theta_0$, $H_1 : \theta = \theta_1$, then among all tests of size  $\alpha$, the likelihood ratio test, defined with test statistic 
\begin{equation*}
    T = \frac{f(x ; \theta_1)}{f(x ; \theta_0)} = \frac{f_1(x)}{f_0(x)}
\end{equation*}
with rejection region $R = \{X : T(X) \geq \kappa_\alpha\}$ where $\alpha \geq P_{H_0}(T(X) \geq \kappa_\alpha)$ is the most powerful test. \\
\textit{Proof.} Let $\tilde R$ be another rejection region with size  $\alpha$. This implies 
\begin{equation*}
    \alpha = \int _R f_0(x) dx = \int_{\tilde R} f_0(x) dx
\end{equation*}
Rewrite 
\begin{align*}
    \int_R f_0(x) dx &= \int_{R \cap \tilde R} f_0(x) dx + \int_{R\cap \tilde R^c} f_0(x) dx\\
    \int_{\tilde R} f(x) dx &= \int _{\tilde R \cap R } f_0(x) dx + \int_{\tilde R \cap R^c} f_0(x) dx
\end{align*}
implies that 
\begin{equation*}
    \int_{R \cap \tilde R}f_0(x) dx = \int_{\tilde R \cap R^c} f_0(x) dx
\end{equation*}
By definition of $R$, for $x \in R$, $f_0(x) \leq \kappa_\alpha f_1(x)$ and in $R^c$, $f_0(x) > \kappa_\alpha f_1(x)$, so 
\begin{equation*}
    \int _{R \cap \tilde R^c} f_1(x) dx \geq \int_{\tilde R \cap R^c} f_1(x) dx
\end{equation*}
Then 
\small
\begin{equation*}
    \int_{R} f_1(x) dx = \int_{R \cap \tilde R^c} f_1(x) dx + \int_{R \cap \tilde R^c} f_1(x) dx \geq \int_{\tilde R \cap R^c} f_1(x) dx + \int_{\tilde R \cap R} f_1(x) dx = \int_{\tilde R} f_1(x) dx
\end{equation*}
\normalsize
hence the test with rejection region $R$ is more powerful than the test with $\tilde R$. 
}

\end{enumerate}
\end{document}

